%==============================================================================
\section{Bộ lọc tương quan cao (High Correlation Filter)}
\label{sec:high-correlation-filter}
%==============================================================================

\begin{dinhnghia}[title={Bộ lọc tương quan cao}]
Kỹ thuật \textbf{lựa chọn đặc trưng} (filter method) dùng để phát hiện và loại các feature \textbf{tương quan mạnh} với nhau. Mục tiêu: giảm số biến dự báo, hạn chế \textbf{đa cộng tuyến} (multicollinearity) --- khi hai hay nhiều biến giải thích gần như cùng một thông tin.
\end{dinhnghia}

\begin{ynghia}[title={Nguyên tắc hoạt động}]
Tính hệ số tương quan giữa từng \textbf{cặp} feature; nếu $|r|$ vượt ngưỡng, loại \textbf{một} trong hai feature của cặp. Feature còn lại dùng huấn luyện mô hình.
\end{ynghia}

\subsection{Thuật toán}

\begin{phuongphap}[title={Các bước triển khai}]
\begin{enumerate}
  \item Tính \textbf{ma trận tương quan} (correlation matrix) trên toàn bộ feature.
  \item Đặt \textbf{ngưỡng} cho hệ số tương quan (ví dụ $0{,}8$).
  \item Tìm các \textbf{cặp feature} có $|r|$ lớn hơn ngưỡng.
  \item Với mỗi cặp tương quan cao, \textbf{xóa một} feature (thường giữ feature có tương quan mạnh hơn với biến mục tiêu, hoặc xóa theo quy tắc cố định như trong ví dụ).
  \item Huấn luyện mô hình trên tập feature còn lại.
\end{enumerate}
\end{phuongphap}

\subsection{Ví dụ Python --- dataset Diabetes}

\begin{dongco}[title={Dữ liệu}]
Dùng bộ \texttt{diabetes} từ OpenML qua \texttt{sklearn.datasets.fetch\_openml}; ngưỡng tương quan $0{,}8$.
\end{dongco}

\begin{vidu}[title={Mã đầy đủ}]
\begin{lstlisting}
import numpy as np
import pandas as pd
from sklearn.datasets import fetch_openml

data = fetch_openml("diabetes", version=1, as_frame=True, parser="auto").frame
X = data.iloc[:, :-1].values
y = data.iloc[:, -1].values

corr_matrix = np.corrcoef(X, rowvar=False)
threshold = 0.8

high_corr_indices = np.where(np.abs(corr_matrix) > threshold)
features_to_remove = set()

for i, j in zip(*high_corr_indices):
    if i != j and (j, i) not in features_to_remove:
        features_to_remove.add((i, j))

features_to_remove = list(features_to_remove)
X_filtered = np.delete(X, [j for i, j in features_to_remove], axis=1)

print('Shape of the filtered dataset:', X_filtered.shape)
\end{lstlisting}
\end{vidu}

\begin{ynghia}[title={Kết quả đầu ra}]
\textbf{Output:}
\begin{lstlisting}
Shape of the filtered dataset: (768, 8)
\end{lstlisting}
Với ngưỡng $0{,}8$, không có cặp feature nào vượt ngưỡng trên bộ diabetes chuẩn hóa từ OpenML, nên \textbf{giữ nguyên 8} biến dự báo và $768$ mẫu.
\end{ynghia}

\subsection{Ưu và nhược điểm}

\begin{tinhchat}[title={Ưu điểm}]
\begin{itemize}
  \item \textbf{Giảm đa cộng tuyến}: loại biến trùng lặp thông tin, ổn định hơn cho hồi quy tuyến tính.
  \item \textbf{Cải thiện hiệu năng mô hình}: ít feature $\Rightarrow$ huấn luyện nhanh hơn, ít nhiễu.
  \item \textbf{Đơn giản hóa mô hình}: dễ giải thích khi số biến giảm.
  \item \textbf{Tiết kiệm tài nguyên tính toán}: ma trận thiết kế nhỏ hơn.
\end{itemize}
\end{tinhchat}

\begin{luuy}[title={Nhược điểm}]
\begin{itemize}
  \item \textbf{Mất thông tin}: feature bị xóa có thể chứa tín hiệu hữu ích riêng.
  \item \textbf{Giả định tuyến tính}: chỉ đo tương quan Pearson; quan hệ phi tuyến có thể bị bỏ sót.
  \item \textbf{Ảnh hưởng biến phụ thuộc}: nếu hai feature tương quan cao đều quan trọng với $y$, xóa một có thể làm giảm chất lượng dự báo.
  \item \textbf{Selection bias}: quy tắc ``xóa feature thứ hai trong cặp'' có thể không tối ưu cho mọi bài toán.
  \item \textbf{Không đảm bảo tập tối ưu}: filter độc lập mô hình, có thể không khớp metric cuối (AUC, F1, \ldots).
\end{itemize}
\end{luuy}
